% Source: 28_Chapter_19_Lifted_In_a_Mystic_Manner_The_Papacy_an.txt
\bookchapter[Lifted in a Mystic Manner]{CHAPTER 19}{LIFTED IN A MYSTIC MANNER}{The Papacy and the Crusader}

\leadin{I}{N PARIS, ON A} small island in the Seine, stands an edifice of weather-beaten stone, the
Gothic Cathedral of Notre Dame. Dedicated to the glory of God and the veneration of “Our Lady,” this
well-known sanctuary reflects a fascinating image of the life and spirit of the Christian Middle
Ages.

Notre Dame was erected between 1163 and 1235 during some of the most epoch-making years of Western
Christianity. While workmen were constructing a flying buttress or carefully fitting a beautiful
window in place, students reclined on the Petit Pont, a bridge that led to the Left Bank, and
exchanged reports of crusader victories in the Holy Land or wrangled over one of the theological
virtues. Some of these students would shortly wear the colors of the Episcopal office, and one of
them would carry the papal office to the zenith of worldly power as Pope Innocent III.

Notre Dame was only one symptom of a veritable fever of church building then infecting Europe.
Between 1170 and 1270 more than five hundred great churches were built in Gothic style in France
alone. Gothic first appeared in the restoration of the abbey church of St. Denis near Paris. The
abbot of St. Denis, Suger, conceived the project and guided the work of rebuilding between 1137 and
1144. He described his reaction to the restoration in terms seldom surpassed in praise of the
Gothic: “I seem to find myself, as it were, in some strange part of the universe which was neither
wholly of the baseness of the earth nor wholly of the serenity of heaven, but by the grace of God I
seemed lifted in a mystic manner from this lower towards the upper sphere.”

Lift. Reach. Aspiration. These are the terms for the age. The Middle Ages were not all dark, as the
hundreds of Gothic cathedrals so readily testify.

What can Christians hope for in human society? If God’s will were done on earth as it is in heaven,
what would earth look like? Believers in every age have asked that question, but no age has reached
for the stars quite like the so-called High Middle Ages.

\section{THY KINGDOM COME}

During the twelfth and thirteenth centuries, the papacy led an admirable attempt to constitute a
perfect society on earth. The church achieved an incomparable power and majesty. Like the Gothic
cathedral the medieval church shot upward into the heavens, calling all below it to the glory of
God. But like the cathedrals, the papacy reached for the impossible and first cracked, then, in
time, crumbled to earth.

It happened often with the churches. Early Gothic builders tried to attain the utmost height
possible with the materials at hand. The Cathedral of Chartres, for example, equals the elevation of
a skyscraper of thirty stories. The Strasbourg Cathedral rises forty stories. Many such efforts
resulted in repeated failures in church construction. One historian of cathedrals has remarked that
the question to put to a tourist guide is, “When did the tower fall down?”

By reaching the greatest height possible, the architect tried to create the illusion of soaring. The
use of flying buttresses eliminated the massive walls of earlier churches and made possible slender
pillars that, in turn, were carved in the form of a cluster of still more slender shafts reaching
up, up into infinity. Even the sculptured forms of saints or angels standing in the niches were
elongated: the neck, the arms, the legs seemed to be soaring heavenward.

That is probably the best way to understand Western Christianity in the twelfth and thirteenth
centuries. The papacy ascended high above European society, mounting upon the fading glory of the
empire. The emergence of unified national states reduced the empire to a shadow of its former
universal power. Emperors continued to call themselves “ever-august Roman Emperors” and continued to
go to Rome for their coronations, but they were actually merely the sovereigns of the cluster of
kingdoms and municipal republics that constituted the Germany of the late Middle Ages.

The papacy, by contrast, built upon the reforms of Pope Gregory VII and emerged as the most powerful
office in Europe. The pope’s government was a truly universal monarchy, steadily becoming completely
centralized. All the bishops swore fealty to the pope, no religious order could be founded without
his authorization, the papal court in Rome heard appeals from all over Christendom, and in every
country legates from Rome watched over the execution of papal orders.

In the hands of a strong leader, the papacy could overshadow all secular monarchs. Such a leader was
Pope Innocent III (1198–1216), a new type of administrator-pope. Unlike Gregory VII and other
earlier reform popes who were monks, Innocent and other great popes of the later twelfth and
thirteenth centuries were trained as canon lawyers, experts in church government. Innocent was like
Gregory VII, however, in holding an exalted view of his office. “The successor of Peter,” he
announced, “is the Vicar of Christ: he has been established as a mediator between God and man, below
God but beyond man; less than God but more than man; who shall judge all and be judged by no one.”

Innocent III told the princes of Europe that the papacy was like the sun, while kings were like the
moon. As the moon received its light from the sun, so kings derived their powers from the pope. The
papacy’s chief weapons in support of this authority were spiritual penalties. Almost everyone
believed in heaven and hell and in the pope’s management of the grace to get to one and avoid the
other.

Thus the pope’s first weapon in bringing peasants and princes to their knees was the threat of
excommunication. He could pronounce their anathema and they would be “set apart” from the church,
deprived of the grace essential for salvation. After some bishop read the solemn sentence of
excommunication, a bell rang as for a funeral, a book was closed, and a candle was extinguished, all
to symbolize the cutting off of the guilty man. If he entered a church during mass he was expelled
or the mass was halted.

While under excommunication, persons could not act as judge, juror, witness, or attorney. They could
not be guardians, executors, or parties to contract. After death they received no Christian burial,
and if, by chance, they were buried in consecrated ground, the church had their bodies disinterred
and destroyed.

The second weapon in the papal arsenal was the interdict. While excommunication was aimed at
individuals, the interdict fell upon whole nations. It suspended all public worship and, with the
exception of baptism and extreme unction, it withdrew the sacraments from the lands of disobedient
rulers. Pope Innocent III successfully applied or threatened the interdict eighty-five times against
uncooperative princes.

So successful was the pontiff in asserting his temporal as well as spiritual supremacy that many
states, both large and small, formally acknowledged the pope as their feudal lord. In England, King
John differed with Innocent over the election of the archbishop of Canterbury, and Innocent placed
England under interdict and excommunicated John. Under attack from his barons, John capitulated to
Innocent by becoming his vassal, receiving England back as a fief, and paying him a sizable annual
tribute.

In France, Innocent forced King Philip Augustus to comply with the church’s moral code by taking
back as his queen the woman he had divorced with the consent of the French bishops.

And within the Holy Roman Empire (Germany), Innocent intervened in a civil war between rival
candidates for the throne, supporting first one, then the other. In the end Innocent secured the
election of his ward, the young Hohenstaufen heir Frederick II, who promised to respect papal rights
and to go on a crusade.

Wielding these spiritual weapons, Innocent and his successors in the papal office during the
thirteenth century led Christianity to its peak of political and cultural influence. We simply
cannot understand the pope’s place in our own times without some understanding of these years.

The soaring dreams—and delusions—of the papacy in this period appear preeminently in the crusades
and in Scholasticism. In one the pope claimed power over the holy cause in history; in the other he
maintained authority over the souls of men even in eternity. We will look at the crusades in this
chapter and examine Scholasticism in the next.

\section{TAKING UP THE CROSS FOR CHRIST}

The crusades reflected the new dynamism in Christianity. Islam had come to birth through the visions
and revelations claimed by its prophet, Muhammad, who was born in AD 570. At the age of twenty-five
he suddenly moved from being a keeper of camels to overseeing the business affairs of a wealthy
widow. She was impressed by his character, and though fifteen years his senior, soon proposed
married to him. Although after she died he took numerous wives, while she was alive she remained his
only wife.

At the age of forty Muhammad began to retire into a cave for extensive periods of contemplation and
meditation. After one of these occasions he told his wife that he had been visited by an angel who
ordered him “to recite,” and from that command came the Koran, which literally means to recite.

At first he was perplexed by it all, but his wife told him this could well be the call of God upon
him. She became his first follower, and then a handful of family members followed suit. Many
resisted his claim to be God’s spokesman, and it was only after various battles, arguments, and wars
that Muhammad was able to gain a vast following in Arabia.

In the first one hundred years of this Arabic movement, numerous capitals, such as Jerusalem,
Damascus, and Cairo, fell to Islam. But the succession to Muhammad after his unexpected death in AD
632 became a serious, divisive issue. He had not named a successor, and an immediate struggle for
the leadership of the movement ensued. One group, the Sunnis, insisted that the first caliph should
be elected. Shiahs or Shiites argued that the successor must come from Muhammad’s own bloodline,
which would have been Ali, his cousin and son-in-law.

Now, in the age of the grand cathedral, crusaders from Western Europe attempted to expel the Muslims
from the Holy Land. Historians have attributed numerous motives to the crusaders beyond their
religious fervor. Did they seek adventure in these strange foreign lands? Would they find personal
gain or spiritual advantage? While few contemporary Christians would defend the idea of the crusades
or its most gross offenses, we must not overlook a simple reality. The Christians sought to counter
Islam’s remarkable military conquest and preserve their geographic strongholds from being overrun.
By the eighth century, one half of all Christians lived under Islamic rule. All the great and
colorful figures of this era were caught up in this consuming cause, from Peter the Hermit, who
enflamed the First Crusade, to the saintly Louis IX, King of France, who inspired the Sixth and
Seventh.

For centuries peaceful pilgrims had been traveling from Europe to worship at the birthplace of
Christ. The rise and spread of Islam in the Near East during the seventh century did not interrupt
this traffic. By the tenth century bishops were organizing mass pilgrimages to the Holy Land; the
largest, which set out from Germany in 1065, included about seven thousand pilgrims.

During the eleventh century, however, Christian pilgrims began to encounter persecution, and when
the Seljuk Turks, new and fanatical converts to Islam, came sweeping and plundering into the Near
East, the situation became especially tense. The Seljuks seized Jerusalem from their fellow Muslims
and then swept north into Asia Minor.

Forces of the Eastern Empire desperately tried to bar the invader, but at the battle of Manzikert
(1071) the Turks captured the eastern emperor and scattered his army. Within a few years Asia Minor,
the chief source of Byzantine revenue and troops, was lost, and the emperor was writing to western
princes and to the pope seeking mercenaries to aid in the rescue of lost territories. In addition,
tales of alleged Turkish mistreatment of Christian pilgrims circulated throughout Europe, and though
there is evidence that these stories were propaganda, rumors themselves were sufficient to inflame
men’s minds.

In 1095, after the Eastern Emperor Alexius I sent out an urgent appeal for help, Pope Urban II
proclaimed the First Crusade to regain the Holy Land. Preaching at the Council of Clermont in
southeastern France, the pope urged Christians to take up the cross and strive for a cause that
promised not merely spiritual rewards but material gain as well:

\begin{blockquote}
For this land which you inhabit~.~.~. is too narrow for your large population; nor does it abound in
wealth; and it furnishes scarcely food enough for its cultivators. Hence it is that you murder and
devour one another~.~.~. enter upon the road to the Holy Sepulchre; wrest that land from the wicked
race, and subject it to yourselves.
\end{blockquote}

As Urban ended his impassioned appeal a roar rose from the multitude: “Deus Vult!” “God wills it!”
So there on the spot Urban declared that Deus Vult! would be the crusader battle cry against the
Muslim enemy.

For seven centuries Christians have tried to forget the crusades, but neither Jew nor Muslim will
allow them to do so. In our liberated generation it is easy to dismiss the whole bloody affair as
insane religious bigotry, forgetting the context in which it occurred.

But the crusaders were human beings, so their motives, like our own, were mixed and often in
conflict. The word crusade itself comes from “taking the cross,” after the example of Christ. Thus
on the way to the Holy Land the crusader wore the cross on his breast. On his journey home, he wore
it on his back.

The crusaders were fully aware of the spiritual rewards Urban promised them, including full
forgiveness of their past sins. And most of them shared a profound personal reverence for the soil
that Christ had trod.

The intensity of this emotion was later caught by Shakespeare in words of that pugnacious English
monarch Henry IV:

\begin{poem}
We are impressed and engag’d to fight~.~.~.\\
To chase those pagans in those holy fields,\\
Over whose acres walked those blessed feet,\\
Which fourteen hundred years ago were nail’d,\\
For our advantage on the bitter cross.
\end{poem}

For Urban and the popes who followed him, the crusades were a new type of war, a Holy War. Augustine
had laid down the principles of a “just war”: It was conducted by the state; its purpose was the
vindication of justice, meaning the defense of life and property; and its code called for respect
for noncombatants, hostages, and prisoners. All these evaporated in the heat of the holy cause.
Urban appealed for crusaders in the name of the church; the purpose of liberating fellow Christians
morphed into the conquest of the infidels in the Holy Land; and this high calling excused somehow
the total disregard for noncombatants and prisoners.

The inception of the crusades ignited horrible attacks against the Jews, and even fellow Christians
were not exempt from rape and plunder. Incredible atrocities befell the Muslim foes. Crusaders sawed
open dead bodies in search of gold, sometimes cooking and eating the flesh—a delicacy they found
“better than spiced peacock,” as one chronicler chose to describe it.

\section{EVENTS OF THE CRUSADES}

From the end of the eleventh century to the end of the thirteenth, Christian Europe, led by the
popes, launched seven major crusades, as well as various small expeditions.

The First Crusade was composed of feudal nobles from France, parts of Germany, and southern Italy,
where the Norman (or Viking) invaders had settled. The armies proceeded overland to Constantinople.
Having expected the help of European mercenaries against the Seljuks, the emperor Alexius Comnenus
was taken aback when confronted by an unruly horde of what Pope Urban himself had called “aforetime
robbers.” The emperor hastily directed the crusaders out of Constantinople to fight the Turks.

The First Crusade was the most successful of the seven; with not more than five thousand knights and
infantry, it overcame the resistance of the Turks, who were no longer united. Above all, it captured
the Holy City, Jerusalem. A contemporary account of the Christian entrance into Jerusalem reads,

\begin{blockquote}
Some of our men~.~.~. cut off the heads of their enemies; others shot them with arrows, so that they
fell from the towers; others tortured them longer by casting them into the flames.~.~.~. It was
necessary to pick one’s way over the bodies of men and horses. But these were small matters compared
to what happened at the Temple of Solomon [where]~.~.~. men rode in blood up to their knees and
bridle reins. Indeed it was a just and splendid judgment of God that this place should be filled
with the blood of the unbelievers, since it had suffered so long from their blasphemies.
\end{blockquote}

\noindent At nightfall the crusaders’ hands were still bloody when they folded them in prayer and
knelt at the Church of the Holy Sepulchre, “sobbing for excess of joy.”

The First Crusade captured a long strip of territory along the eastern coast of the Mediterranean
and created the feudal Latin kingdom of Jerusalem. It survived until 1291 when its last remnant fell
to the Muslims.

When the kingdom of Jerusalem faced its first crisis, tottering on the brink of destruction in 1147,
Bernard, the powerful mystic of Clairvaux, called for the Second Crusade. Despite Bernard’s rhetoric
and the presence of royalty, the crusade achieved nothing. After two years it simply melted away.

The original frenzy had clearly cooled, and the signs of corruption of the holy cause were apparent.
The popes needed money to meet such obligations as providing legates for the new Christian lands in
the East. So they turned spiritual benefits into moneymaking advantage.

In the Middle Ages sins to be forgiven had to be confessed. Upon hearing a confession the priest not
only pronounced the penitent’s guilt for sin forgiven (thanks to the merit of Christ) but also
required a penalty or satisfaction—some “act of penance”—as a mark of the penitent’s sincerity. If
the penitent died before performing this penance, purgatory afforded him an opportunity in the life
after death. Such penance, whether in this life or in purgatory, was called “temporal” punishment.

For years the church had claimed the power to remit part of this temporal punishment, but no
complete remission had been granted until Urban II at Clermont offered total remission or
“indulgence” for crusaders who headed for Jerusalem “out of pure devotion.”

It was only a slight step farther to confer like benefits upon those who were unable to go on a
crusade but who contributed to the cause. A man could virtually buy a substitute. Thus the
possibilities for fund-raising opened in all directions, including construction of a hospital here
or a cathedral there.

In 1187 Saladin, the Sultan of Egypt and Syria, brought fresh and vigorous leadership to the
Muslims. When Jerusalem fell to the infidels, Christians with some reluctance responded to the cry
for the Third Crusade (1189). Its leaders were three of the most famous medieval kings: Frederick
Barbarossa of Germany, Richard the Lion-Hearted of England, and Philip Augustus of France. Frederick
was drowned in Asia Minor; and, after many quarrels with Richard, Philip returned home. Saladin and
Richard remained the chief protagonists.

To keep the Muslims united, Saladin proclaimed a jihad, or holy war, against the Christians, but he
remained a patient statesman and chivalrous warrior. “Abstain from the shedding of blood,” he once
said, “for blood that is spilt never slumbers.” His commonsense approach to a settlement was evident
when he proposed that Richard should marry his sister and be given Palestine as a wedding present, a
proposal that shocked the Europeans. Richard and Saladin finally agreed to a three-year truce and
free access to Jerusalem for Christian pilgrims. Since Saladin would have granted this concession at
any time, the truce scarcely compensated for the cost of such an expensive crusade.

The Fourth Crusade revealed the harsh realities of papal aspirations in the East and the degradation
of the crusading zeal. Upon ascending the papal throne in 1198, Innocent III threw himself into the
effort of reviving the crusading spirit, but the few knights who answered his call were unable to
meet the outrageous shipping charges demanded by the Venetians. To secure the costs of shipping, the
Venetians persuaded the crusaders to capture the Christian town of Zara on the Adriatic coast.

The town had been a nuisance to Venetian vessels for years. So in 1202 the crusaders sacked Zara.
Innocent complained that Satan was behind the whole rapacious affair and proceeded to excommunicate
all who participated. The Venetians, however, followed their advantage by pressuring the crusaders
to attack Constantinople itself. Torn by warring parties within, Constantinople fell to the
crusaders and, after ravaging the city, the crusaders set up the Latin Empire of Constantinople in
1204 and promptly forgot about “the deliverance of the Holy Land.”

Upon hearing of the conquest Innocent wrote in rage, “You have spared nothing that is sacred,
neither age nor sex. You have given yourselves up to prostitution, to adultery and to debauchery in
the face of all the world.” Never one to pass a political advantage, however, Innocent appointed an
archbishop of Constantinople to serve Roman interests.

The Latin Empire in Constantinople lasted until 1261, but the ancient city never fully recovered.
The conquest widened the schism between the Greek and the Latin churches and hastened the fall of
the city in 1453 to the Turks.

Other crusaders marched and meandered east during these years but none of the holy efforts postponed
the predictable, the return of the Holy Land to Muslim control. The era of crusades ended in 1291
when Acre, the last stronghold of the Christians in the land where Jesus walked, fell to those who
denied his deity.

\section{RESULTS OF THE HOLY CAUSE}

The long-range results of two centuries of crusading zeal are not impressive. If the primary
purposes of the crusades were to win the Holy Land, to check the advance of Islam, and to heal the
schism between the Eastern and Western churches, then the crusades failed miserably.

The crusades created three semimonastic military orders: the Templars (or Knights of the Temple)
whose first headquarters were on the site of the old Temple of Jerusalem; the Hospitalers (or
Knights of St. John of Jerusalem) who were founded originally to care for the sick and wounded; and
the Teutonic Knights, exclusively a German order. Combining monasticism and militarism, these orders
had as their aims the protection of all pilgrims and perpetual war against the Muslims. These men of
the cross could put five hundred armed knights into the field, and their great castles guarded the
roads and passes against Muslim attack. For two centuries the Templars in white robes decorated with
a red cross, the Hospitalers in black robes with a white cross, and the Teutonic Knights in white
robes with a black cross were common sights in the Crusader States.

\figtitle{The Age of Crusades}

\figtitle{The High Middle Ages}

Perhaps the most significant result was the added splendor the crusades brought to the papacy. Not
only did a pope, Urban II, launch the First Crusade, but the popes throughout the period were the
primary inspiration for fresh expeditions. They, not the emperors, strove to unite Christendom
against Islam. The new military orders and the new bishops in the Holy Land and those in
Constantinople for a time were under papal protection and service. The Holy War was the papacy
reaching for universal sovereignty, one united Church, West and East.

But like the Gothic architects who built their cathedrals higher and higher until the towers cracked
and then crashed to earth, the popes tried the impossible. Christian Europe had no need for Syria or
Jerusalem. She took them in a fit of enthusiasm and had no power to retain them. When trade and
towns turned rulers and people to new interests at home, the popes remained faithful to the old
ideal, Christian control of the Holy Land. It was their constant preoccupation.

Unfortunately the popes never held two basic truths that we must never forget: Christianity’s
highest satisfactions are not guaranteed by possession of special places, and the sword is never
God’s way to extend Christ’s church. This fault assured the religious collapse of the whole
structure.

\readinghead{Suggestions for Further Reading}
\begin{readinglist}
\entry Barraclough, Geoffrey. The Medieval Papacy. New York: Harcourt, Brace \& World, 1968.
\entry Madden, Thomas. A Concise History of the Crusades. New York: Rowman \& Littlefield, 1999.
\entry Riley-Smith, Jonathan. The Crusades: A Short History. New Haven, CT: Yale University Press,
  1987.
\entry *Stark, Rodney. The Battalions: The Case for the Crusades. New York: Harper One, 2012.
\end{readinglist}
